% Class-compatibility preamble retained for a reader excerpt of
%   Sahar Diskin, Anna Geisler, "Perfect Matching in Product Graphs and in their
%   Random Subgraphs", arXiv:2404.14020v4, CC BY 4.0.
%
% The source pins its own publisher class (scrartcl); the delivery wrapper accepts
% only the classes article and amsart, so this file supplies the amsart baseline and
% then copies the author's own declaration block mechanically and unchanged
% (pinned-original-source.tex, source lines 30-57: the author's newtheorem
% declarations and the author's macro declarations).  The author's own package
% list of source lines 4-21 is reused only for the entries the retained text needs
% (enumitem, amsmath, amssymb, amsthm, mathtools, xcolor, geometry) and the
% cross-reference packages already needed by the author's own \Cref/\eqref usage
% (hyperref, cleveref); tikz, caption, subcaption, float, graphicx, wrapfig and
% abstract serve figures and the title block that this excerpt does not carry.
% No mathematical macro, statement, symbol or assumption is introduced, changed
% or removed by this file.
\documentclass[10pt,a4paper]{amsart}
\usepackage[shortlabels]{enumitem}
\usepackage{amsmath}
\usepackage{amssymb}
\usepackage{amsthm}
\usepackage{mathtools}
\usepackage{xcolor}
\usepackage[margin=1in]{geometry}
\usepackage{hyperref}
\usepackage{cleveref}
\hypersetup{
  colorlinks=true,
  linkcolor=blue,
  filecolor=magenta,
  urlcolor=cyan
}

\newtheorem{theorem}{Theorem}
\newtheorem{thm}{Theorem}[section]
\newtheorem{corollary}[thm]{Corollary}
\newtheorem{lemma}[thm]{Lemma}
\newtheorem{proposition}[thm]{Proposition}
\newtheorem{claim}[thm]{Claim}
\newtheorem{remark}[thm]{Remark}
\newtheorem{question}[thm]{Question}

\newcommand{\whp}{whp}
\newcommand{\defi}{\operatorname{def}}
\newcommand{\odd}{\operatorname{odd}}
\newcommand{\iso}{\operatorname{iso}}
\newcommand{\cN}{\mathcal{N}}
\newcommand{\cO}{\mathcal{O}}
\newcommand{\cE}{\mathcal{E}}
\newcommand{\cK}{\mathcal{K}}
\newcommand{\cP}{\mathcal{P}}
\newcommand{\cQ}{\mathcal{Q}}

\newcommand{\footremember}[2]{%
  \footnote{#2}%
  \newcounter{#1}%
  \setcounter{#1}{\value{footnote}}%
}

\newcommand{\ag}[1]{\textcolor{purple}{(Anna: #1)}}
\newcommand{\sd}[1]{\textcolor{red}{(Sahar: #1)}}
\title{In a high-dimensional product graph, two isolated vertices of the random subgraph are typically at distance at least three}
\author{Sahar Diskin, Anna Geisler}
\date{arxiv:2404.14020v4; CC BY 4.0}
\begin{document}
\maketitle
\noindent\textit{Editorial note.} All delivered mathematical text is the author's own, copied from the pinned author source sources/combinatorics-fl04-vspace-1-5cm-perfect-matching-in-p/pinned-original-source.tex; no statement or proof step was rewritten, reordered or removed.\par\medskip
\noindent\textit{Editorial note.} Preamble provenance: as for the other cards of this source, the source pins the publisher class scrartcl, which the delivery wrapper does not accept (supported wrapper classes are article and amsart), so this card uses sources/combinatorics-fl04-vspace-1-5cm-perfect-matching-in-p/retained-preamble.tex, which supplies the amsart baseline and then carries the author's own newtheorem and macro declarations (pinned source lines 30-57) verbatim. No macro was invented and no mathematical symbol was changed.\par\medskip
\noindent\textit{Editorial note.} Bibliography provenance: the counted statement and its proof contain no citation, so this card carries no bibliography. The source itself has no thebibliography environment and ends with bibliographystyle abbrv plus bibliography perc (author source lines 1691-1692); no entry was reconstructed or invented.\par\medskip
\noindent\textit{Editorial note.} The displayed range for p that the counted proof invokes, source label eq:mu-range, is delivered verbatim as a context span so that the last step of the proof refers to a formula present in the assembled material; it is the author's own text and is not a reconstruction.\par\medskip
\noindent\textit{Editorial note.} The constant alpha\_M and the parameters n, d and p that the proof uses are set up in the author's model description earlier in the paper; the editorial prerequisite list on the first page names them and gives the source location, and no author text was moved into the counted spans to replace them.\par\medskip
\noindent\textit{Allowed standard prerequisites (editorial).}\begin{itemize}
\item The parameters of the model: n the order and d the dimension-regularity parameter of the product graph, p the edge probability, and the constant alpha\_M controlling the one-edge overlap factor. Used inside the counted proof at source lines 494-508; the model and the relation d = Theta\_M(log n) are set up earlier in the paper (around source lines 201-240). Gives the two counts of ordered pairs (nd adjacent pairs and nd squared pairs at distance two) and the constant alpha\_M in the bound (1-p)\textasciicircum{}(2d-1) <= alpha\_M (1-p)\textasciicircum{}(2d).
\end{itemize}\medskip
\begin{lemma}\label{lem:isolates-far}
Whp, every two isolated vertices of $G_p$ are at distance at least three in $G$.
\end{lemma}

\begin{equation}\label{eq:mu-range}
  d^{-1/50}\le n(1-p)^d\le d^{1/50}.
\end{equation}

\begin{proof}
There are at most $nd$ ordered adjacent pairs and at most $nd^2$ ordered pairs at distance two.
If distinct vertices $u,v$ are not adjacent, the sets of edges incident with them are disjoint, and hence
\[
  \mathbb{P}(u,v\text{ are isolated in }G_p)=(1-p)^{2d}.
\]
If $u$ and $v$ are adjacent, their incident edge sets overlap in one edge, so this probability is
\[
  (1-p)^{2d-1}\le \alpha_M(1-p)^{2d},
\]
for some constant $\alpha_M$.
The expected number of ordered pairs of isolated vertices at distance at most two is therefore at most
\[
  \alpha_M nd^2(1-p)^{2d}   \le \frac{\alpha_M d^{2+1/25}}{n}=o(1),
\]
where we used that $p$ satisfies \eqref{eq:mu-range}.
We conclude with Markov's inequality.
\end{proof}

\end{document}
